\qquad\brandName realizó pruebas bajo un enfoque de “\typePentest” desde \initialDate hasta \finalDate, sin credenciales ni conocimiento previo del entorno interno de \enterpriseName, con el objetivo de identificar debilidades desconocidas. Las pruebas se realizaron desde un punto de vista no invasivo con el objetivo de descubrir la mayor cantidad posible de configuraciones erróneas y vulnerabilidades. Las pruebas se llevaron a cabo de forma remota a través de un host que fue aprovisionado específicamente para esta evaluación. Cada debilidad identificada fue documentada e investigada manualmente para determinar las posibilidades de explotación y el potencial de escalamiento. 

\brandName buscó demostrar el impacto total de cada vulnerabilidad, incluyendo el compromiso del dominio interno. Si \brandName lograba obtener un punto de apoyo en la red interna, \enterpriseName permitió realizar pruebas adicionales, incluyendo movimiento lateral y escalamiento de privilegios horizontales/verticales, para demostrar el impacto de un compromiso de la red interna.
